\subsection{一些预等变膜与等变膜}\label{sec:18e}
现在重新考察\secref{18a}中引入的 $\widetilde E$ 中两类 Lagrangian 子流形，目的是将它们变成等变 Fukaya 范畴 $\F(\widetilde E)^{\mathrm{naive}}$ 的对象。回顾，为定义该范畴，必须对 $\Z/2$ 作用施加一些技术条件，见\stref{假设}{14.3}。在当前情形中，这些条件显然满足：$\iota_{\widetilde E}$ 的不动点集为 $\{y=0\}\cong M$，其法丛有自然平凡化，由 $D\widetilde\pi|M:T\widetilde E|M\to T\widetilde S_0=\CC$ 给出。

\begin{description}[leftmargin=2em]
\item[(U) 型。]\mbox{}\par 设 $L^\#=(L,\alpha^\#,P^\#)$ 为 $E$ 中的正合 Lagrangian 膜，包含于 $E^{\mathrm{shrink}}$ 的内部。设 $\widetilde L$ 为 $L$ 在 $\widetilde E$ 中的原像。由\eref{18.15}，$L$ 与 $\widetilde L$ 的平方相位映射满足
\[
\alpha_{\widetilde E}(T\widetilde L_{(y,x)})=\frac{|y|^2}{y^2}\alpha_E(TL_x).
\]
回顾，$\widetilde L=\widetilde L_+\cup\widetilde L_-$，其中连通分支 $\widetilde L_\pm$ 位于子集 $\pm\operatorname{im}(y)>0$ 内。由于 $\operatorname{im}(y)$ 始终非零，可以取辐角函数 $\arg:\CC\setminus\R^+\to(0,2\pi)$，令
\begin{equation}\label{eq:18.18}
\widetilde\alpha^\#(y,x)=\alpha^\#(x)-\frac1{2\pi}\arg(y^2),
\end{equation}
从而定义 $\widetilde L$ 的分次。限制后，得到 $\widetilde L_\pm$ 的分次 $\widetilde\alpha^\#_\pm$。至于 Pin 结构，我们约定为 $\widetilde L_+$ 赋予给定 $P^\#$ 的拉回，为 $\widetilde L_-$ 赋予 $P^\#\otimes\lambda^{\mathrm{top}}(TL)$ 的拉回。将它们记为 $\widetilde P_\pm^\#\to\widetilde L_\pm$，将其并记为 $\widetilde P^\#\to\widetilde L$。

据此，容易将 $\widetilde L_+$ 赋予\secref{14a}意义下的预等变膜结构。其主体为 $\widetilde L_+^\#$，影子为 $\widetilde L_-^\#$，同构 $\widetilde f_-:\widetilde P_-^\#\to\iota_{\widetilde E}^*(\widetilde P_+^\#\otimes\lambda^{\mathrm{top}}(T\widetilde L_+))$ 为恒等同构，并调整 $\widetilde f_+$ 使\eref{14.1}成立。交换 $+$ 与 $-$ 后，相同程序给出 $\widetilde L_-$ 的预等变膜结构；取二者的不交并，$\widetilde L$ 就成为等变膜。
\end{description}
\SourcePage{285}
已经可以得出一些初步结论。设 $L_k^\#=(L_k,\alpha_k^\#,P_k^\#)$，$k=0,1$，为 $E$ 中上述类型的两个膜。为该子流形对选择 Floer 数据 $(H,J)$ 时，可以要求在 $E^{\mathrm{shrink}}$ 之外 $H$ 为零、$J$ 与 $I_E$ 一致。这不妨碍模空间 $\M_Z(y_0,y_1)$ 的正则性，因为在 $L_0\cup L_1$ 的邻域内仍可自由选择。对于原像对 $(\widetilde L_{0,+},\widetilde L_{1,+})$，取提升 $(H,J)$ 得到的 Floer 数据 $(\widetilde H,\widetilde J)$；由于所取 Floer 数据在分歧集附近平凡，这没有问题。每个 $\widetilde y_+\in\C(\widetilde L_{0,+},\widetilde L_{1,+})$ 始终落在 $\widetilde E^{\mathrm{shrink}}_+$ 内，并投影到某个 $y\in\C(L_0,L_1)$，从而建立了两个集合之间的双射。该对应保持绝对指标，并在典范同构意义下保持定向空间：
\begin{equation}\label{eq:18.19}
i(\widetilde y_+)=i(y),\qquad o(\widetilde y_+)\cong o(y).
\end{equation}
这里有一个小问题，可作如下处理：虽然 $\eta_{\widetilde E}^2|\widetilde E^{\mathrm{shrink}}_+$ 不等于 $\eta_E^2$ 的拉回，但可选一个将 $\CC\setminus\R^+$ 变形收缩到点 $-1$ 的映射，并将其作用于\eref{18.15}中的 $|y|^2/y^2$ 项，从而将一个形式变形为另一个，允许相差一个整体符号。这个二次复体积形式的同伦诱导分次的同伦；事实上，它将 $\widetilde L_k$ 的分次 $\widetilde\alpha_{k,+}^\#$ 变形为 $\alpha_k^\#-1/2$ 的拉回。Pin 结构当然本来就一致，因此\eref{18.19}容易成立。

对每个 $u\in\M_Z(y_0,y_1)$，映射 $v=\pi\circ u:Z\to S$ 在 $\Omega=v^{-1}(S\setminus S^{\mathrm{shrink}})$ 上全纯，故可对 $\overline\Omega$ 上的 $|v|$ 应用最大值原理。由于 $\Omega$ 相对紧且不与 $\partial Z$ 相交，结论是 $\Omega=\varnothing$，因此 $u$ 仍落在 $E^{\mathrm{shrink}}$ 内。这意味着，投影在这样的 $u$ 与元素 $\widetilde u_+\in\M_Z(\widetilde y_{0,+},\widetilde y_{1,+})$ 之间建立双射。该双射还将相应线性化算子等同，故 $(\widetilde H,\widetilde J)$ 仍然正则。最后，这一等同与用来定义 Floer 微分 $\partial$ 的贡献\eref{12.17}的粘合过程相容，从而定义了 Floer 上链复形的同构。

对原像的 $-$ 部分，相同论证也适用；唯一区别是 Pin 结构被 $\lambda^{\mathrm{top}}(TL_k)\cong\xi_E|L_k\cong\xi_{\widetilde E}|\widetilde L_{k,-}$ 扭曲。但这一扭曲同时作用于涉及的两个 Lagrangian 子流形，故 Floer 复形不受影响。总之，存在典范同构
\begin{equation}\label{eq:18.20}
\begin{aligned}
CF^*(\widetilde L_0^\#,\widetilde L_1^\#)
&\cong CF^*(\widetilde L_{0,+}^\#,\widetilde L_{1,+}^\#)\oplus CF^*(\widetilde L_{0,-}^\#,\widetilde L_{1,-}^\#)\\
&\cong CF^*(L_0^\#,L_1^\#)\otimes_{\K}\K[+,-],
\end{aligned}
\end{equation}
其中 $\K[+,-]$ 指具有两个生成元 $+$ 与 $-$ 的二维向量空间。\znote{277}{两个生成元 $+$ 与 $-$ 都在 $0$ 次：它们不是次数，也不是 §12 的 $[m]$ 平移记号，更不是群环。之所以能写成张量积，是因为\eref{18.20}的两个直和项次数分布完全相同（$\widetilde\alpha_{\pm}^\#$ 由同一个 $\widetilde\alpha^\#$ 限制而来）。与它配套的是下一句：$\widetilde L^\#$ 的等变性给出交换这两个生成元的对合，取不变部分便回到 $CF^*(L_0^\#,L_1^\#)$；命题 18.13 的证明正是靠这一点，把 $\widetilde{\A}^{\mathrm{inv}}$ 与 $E$ 中由较小膜构成的全子范畴等同起来。[来源：原书 §18e、命题 18.13（印刷页 277、282）]}由于 $\widetilde L_k^\#$ 的等变性，Floer 复形带有典范对合；关于\eref{18.20}，它恰为交换生成元 $+$ 与 $-$ 的映射。更一般地，给定任意同类 Lagrangian 膜的 $(d+1)$ 元组，可以选择具有与上述 Floer 数据相同性质的扰动数据。这样，扰动全纯 $(d+1)$ 边形也被限制在 $E^{\mathrm{shrink}}$ 内，意味着\eref{18.20}与 $\Ai$ 结构相容。
\SourcePage{286}
\StatementNumber{8}\begin{note}\label{st:18.8}
假设所有二次复体积形式都存在连续平方根，使各膜自然带有定向，见\stref{笔记}{11.18}。更明确地说，某点处切空间 $T\widetilde L_{(y,x)}$ 的基 $(v_1,\ldots,v_{n+1})$ 为正定向，当且仅当
\[
\exp(\pi i\widetilde\alpha^\#(y,x))=
\frac{\eta_{\widetilde E}(v_1\wedge\cdots\wedge v_{n+1})}{|\eta_{\widetilde E}(v_1\wedge\cdots\wedge v_{n+1})|}.
\]
从 $(y,x)$ 换到 $(-y,x)$ 时，$\eta_{\widetilde E}=\beta_{\widetilde E}^*\eta_E/2y$ 的符号改变，而按定义，$\widetilde\alpha^\#$ 保持不变。因此，两个分支 $\widetilde L_+$ 与 $\widetilde L_-$ 带有相反定向。
\end{note}
